% =====================================================================
%  PRIMER PART II -- FINDING THE QUESTIONS
%  Comes AFTER Part I (background), BEFORE Part III (results/proofs).
%  Purpose: take a reader who has the background and teach them to
%  GENERATE and CHOOSE the project's research questions themselves,
%  WITHOUT revealing the answers (those live in Part III).
%
%  This is a FRAGMENT: no \documentclass, no \begin{document}.
%  The assembler supplies the preamble, which already defines:
%  definition, theorem, lemma, proposition, corollary, example
%  (printed as "Worked example"), remark, problem, and the macros
%  \R \Z \E \sign \Lip \one \conv \dist \Orb \fdc \Vinv \Vac \Sep
%  \Piinv etc.  "Find-the-question" prompts use the {problem}
%  environment plus \paragraph.
%
%  Citation keys used appear in the comment block at the END of file.
%  NEW keys (not yet in primer.bib) are flagged there.
% =====================================================================

% NOTE: the assembler (primer_v2.tex) already emits the \part* heading and its
% table-of-contents line for Part II, so this fragment must NOT repeat them.

Part~I handed you a toolbox: a way to measure how robust a classifier is, the Lipschitz--margin
inequality (a margin divided by a Lipschitz constant lower-bounds the certified radius) that turns
a smoothness budget and a margin into a certified radius,
the Fourier/projector machinery that says what a symmetry preserves, and the implicit-bias and
NTK lenses that say where gradient descent actually lands. This part does something different.
It does \emph{not} give you results. It puts you in the chair of the person who had to decide
\emph{which questions to ask in the first place}, hands you the same situation they faced, and asks
you to form your own hypotheses before Part~III tells you what happened.

That is a deliberate choice. The fastest way to misread a research paper is to learn its answers
and then reverse-engineer questions that make those answers look inevitable. The answers almost
never were inevitable. They were one defensible bet among several, taken at a fork where reasonable
people would have gone different ways. If you can stand at that fork yourself and feel the pull of
each branch, you will read Part~III as a sequence of decisions rather than a sequence of
revelations, and you will be able to do the same thing on a problem nobody has solved yet.

So: read this part with a pencil. Every research question below is presented as a genuine open
problem with two to four candidate directions and their trade-offs. At each one you will be asked to
commit, on paper, to the direction \emph{you} would take and why, before turning to Part~III. The
goal is not to guess the ``right'' answer. It is to develop the instinct that lets you generate good
questions and triage them.

%======================================================================
\section{The situation: a clean story that broke}
\label{sec:p2-situation}

Research questions do not come from the sky. They come from a concrete irritant: something a
previous effort could measure but not explain. Here is the irritant, told as it actually unfolded,
because the shape of the chain is itself a lesson in where questions live.

\subsection{Act I --- a tidy hypothesis, reproduced}
\label{sub:p2-act1}

The starting point was a clean and quotable claim from the literature: \emph{shift invariance can
reduce adversarial robustness} \citep{ge2021shift}. The mechanism, for the linear case, is exactly
the kind of thing Part~I prepared you to read. A classifier that is invariant to all circular
shifts of its input must give the same answer to a signal and to every shift of it; for a
\emph{linear} classifier this forces the weight vector into the fixed subspace of the shift group,
which (Part~I, the Fourier section) is spanned by the constant vector $\one$. The model can then see
the data only through its mean, the DC (zero-frequency) component, and discards all the rest. On a
synthetic ``one black pixel versus one white pixel'' problem this collapses the achievable max-margin
from a constant down to $2/\sqrt{d}$ in the image dimension $d$ \citep{ge2021shift}. More invariance,
smaller margin, less robustness. A satisfying monotone story.

A team reproduced it. On small image families (MNIST, Fashion-MNIST) trained from a shared template,
they measured a \emph{shift-consistency} score --- the fraction of inputs whose predicted label is
unchanged when the input is shifted, an empirical stand-in for ``how invariant did this model turn
out to be'' --- and they measured adversarial robustness with a projected-gradient (PGD) attack.
Across the families the headline ordering held: \emph{higher consistency went with lower robustness},
just as predicted.

\subsection{Act II --- the deviations it could not name}
\label{sub:p2-act2}

Then the tidy story sprang leaks, and the leaks are the interesting part.

First, the ordering was only \emph{mostly} right. Specific architectures sat in the wrong place:
a couple of convolutional variants were more robust than their consistency score said they should
be, or less. The reproduction could record the deviation but had no second axis to explain it.
The hypothesis ``consistency orders robustness'' was being falsified at the margins by its own data,
and nobody could say what the better axis was.

Second, to turn the picture into something actionable --- ``give me a model that is \emph{both}
shift-invariant and robust'' --- the team drew a two-dimensional scatter (consistency on one axis,
robust accuracy on the other), split it into four quadrants using a fully-connected baseline as the
origin, and aimed for the desirable corner. It was a useful screening picture. But it was a picture,
not a principle: nothing in it said \emph{why} a given architecture lands in a given quadrant, or
whether the quadrants would survive a stronger attack, or whether they would persist as models grow.
The diagnostic worked and was unjustified at the same time, which is the most uncomfortable place a
result can sit.

Third, and most honestly, the write-up's own appendix for theory read, in effect, ``to do.''
Everything was empirical, on tiny images, with a single attack family. There was no certificate, no
worst-case attack, no measurement of the very margin whose collapse the base theory was about.

\paragraph{The lesson of Act II.} Notice where the questions are being born. They are not born in the
parts that worked. They are born at the \emph{deviations} (the models in the wrong quadrant), at the
\emph{unjustified tool} (the quadrant diagram), and at the \emph{explicit hole} (the empty theory
appendix). A reproduction that merely confirms a result generates no questions. A reproduction whose
seams show generates all of them. Train yourself to read for the seams.

\subsection{Act III --- the literature contradicts itself}
\label{sub:p2-act3}

If the base story were the whole story, the fix would just be ``measure more carefully.'' It is not,
because the supervising literature does not agree with itself. Four papers, read precisely, point in
different directions while using the same word, ``invariance.''

\begin{itemize}[leftmargin=1.4em]
  \item \textbf{Invariance hurts --- conditionally.} \citet{ge2021shift} is the base result above.
        But read past the headline and Ge themselves add a crucial nuance: on a dataset whose
        classes live in orthogonal frequencies, where shifting does \emph{not} inflate the effective
        dimension, the convolutional network is \emph{more} robust than the fully-connected one. Their
        own conclusion is that invariance hurts only when it shrinks the margin, and that
        ``the linear margin predicts robustness better than the dimension of the data''
        \citep{ge2021shift}. The base paper's own fine print already suspects the headline axis is
        wrong.
  \item \textbf{Invariance trades off --- provably.} \citet{kamath2020} show empirically that training
        for invariance to larger rotations shrinks the average distance to the decision boundary and
        hurts $\ell_\infty$ robustness; \citet{kamath2021} prove a trade-off on a synthetic
        distribution built from a cyclic error-correcting code, and offer a curriculum that buys a
        better Pareto point without beating the single-objective adversarial baseline. Their
        ``rate of invariance'' $\Pr(f(TX)=f(X))$ is structurally the same quantity as the
        consistency score from Act~I --- yet neither Kamath paper cites Ge. Two independent lines,
        one seam between them.
  \item \textbf{A specific architecture choice causes vulnerability.} \citet{galloway2019} argue that
        batch normalization, by fixing each neuron's first two moments, destroys input-space
        distances and drops PGD robustness by amounts (roughly $8$--$20\%$ across datasets)
        comparable to the architecture effects Act~I reported. This is not a contradiction so much as
        a warning: any cross-architecture robustness comparison is confoundable as a normalization
        effect unless normalization is controlled.
  \item \textbf{Invariance helps.} The anti-aliasing line says the opposite of Ge. Inserting a
        low-pass blur before downsampling restores shift-equivariance, and \citet{zhang2019blurpool}
        reports it \emph{improves} shift-consistency and clean accuracy and average-case corruption
        robustness at once. Crucially, Zhang makes no worst-case $\ell_p$ adversarial claim; that is
        left as future work. The natural ``invariance helps'' counter-voice is about
        \emph{average-case} robustness, not worst-case.
\end{itemize}

State the contradiction plainly. Same word, opposite verdicts: one camp says (conditional)
invariance hurts worst-case $\ell_p$ robustness; another reports it helps; a third proves a
trade-off on a cyclic-shift construction; and two of the three do not cite the base result. A
contradiction in the literature is not a nuisance to be smoothed over in related work. It is a
flashing sign that some shared assumption is wrong, and the assumption is almost always a word doing
too much work. Here the overloaded word is ``invariance.''

\subsection{Act IV --- the gap the wider field leaves open}
\label{sub:p2-act4}

Widen the search and the contradiction sharpens into a precise hole. There is a theory camp that
shows invariance or implicit bias can hurt robustness --- but never about shift. \citet{tramer2020}
prove that defending against $\ell_p$ sensitivity forces \emph{excessive invariance}, opening a
distinct invariance-based attack that changes the true label while the prediction stays put; but
their ``invariance'' is over-invariance to an $\ell_p$ ball, not a group symmetry.
\citet{frei2023} prove that gradient descent's max-margin implicit bias on two-layer ReLU networks
reaches solutions that generalize yet are \emph{provably non-robust}, even when robust solutions
fitting the same data exist; \citet{li2025feature} pin the mechanism on feature averaging and
recover robustness with finer supervision; \citet{minvidal2024} flip the bias from non-robust to
robust by swapping the activation. All three are about the optimizer's bias on synthetic, near-
orthogonal clusters \emph{with no group symmetry}, under $\ell_2$. Their lever on robustness is
activation or supervision, not translation.

And there is an empirical/linear camp that says symmetry helps --- but not in a way that closes the
gap. \citet{sahagokhale2024} state verbatim that ``better shift invariance is generally correlated
with better adversarial robustness,'' and show the effect is architectural rather than a capacity
artifact --- but it is a pure correlation, with PGD/FGSM only (so possibly gradient-masked, see
Part~I), one backbone, no causal isolation, and no theorem linking shift to any radius.
\citet{wang2025bridging} report that rotation/scale equivariance helps without adversarial training,
but that is the \emph{wrong group} (CNNs are already translation-equivariant), it is enforced rather
than emergent, and again no $\ell_p$ theorem. \citet{chenzhu2023} and \citet{lawrence2021} prove
that linear equivariant networks reach a wider margin or a characterized implicit bias --- but they
are \emph{linear only} and stop at margin/generalization, never an $\ell_p$ radius.

\paragraph{The gap, in one sentence.} Across every paper on both sides, the same thing is missing:
\emph{no result says what discriminative signal survives a given translation-invariance, turns that
into an $\ell_p$ robustness certificate, and validates it under a standardized strong attack at
matched capacity.} The ``helps'' side is empirical, or linear, or the wrong group. The ``hurts''
side is about $\ell_p$-ball over-invariance or optimizer bias on symmetry-free clusters. The seam
between them is exactly where the Act~I reproduction's monotone story tore.

This is the situation. A reproduced result whose deviations have no second axis; a working diagnostic
with no principle; a four-paper contradiction; and a field-wide gap with a precise shape. Everything
that follows in this part is a question you could have asked while standing here. Now you get to ask
them.

%======================================================================
\section{Five open problems, and the forks inside them}
\label{sec:p2-problems}

For each question I give you the \emph{situation} (what you know and what is unsettled), the
\emph{candidate directions} you could take with their honest trade-offs, and a prompt asking you to
commit before reading on. I will not name the quantity, the theorem, or the experiment that resolves
it. Resist the urge to flip ahead. The value is entirely in deciding under uncertainty.

A note on the format. ``Direction'' means a research strategy, not a guaranteed win. Each has a cost.
Part of taste is reading those costs accurately, because a direction that cannot be wrong is usually
a direction that cannot teach you anything either.

%----------------------------------------------------------------------
\subsection{Q1 --- What single quantity actually orders robustness?}
\label{sub:p2-q1}

\paragraph{Situation.} Act~I gave you a metric, shift-consistency, that orders robustness
\emph{most} of the time and is falsified at the margins by its own outliers. Act~III gave you the
base paper quietly conceding that ``the linear margin predicts robustness better than the dimension
of the data.'' You suspect consistency is a symptom, not a cause. The deeper question is whether a
\emph{single measurable scalar} even exists that you can attach to a trained model and that ranks
models by their robust radius --- and if one does, which family of quantities it belongs to. You are
genuinely choosing among candidate \emph{families} here, not polishing a known winner: a
\emph{consistency} score, a bare \emph{margin}, a \emph{sensitivity} measure (a Lipschitz or
gradient-norm proxy), or some \emph{ratio} of a margin to a sensitivity. Each is a different bet
about what governs the robust radius, and when your chosen scalar disagrees with consistency, you
want it to be the one that is \emph{right}.

\paragraph{Candidate directions.}
\begin{itemize}[leftmargin=1.4em]
  \item \emph{Refine consistency.} Keep the consistency axis but measure it better --- average over
        more shifts, weight by frequency, use a soft (probability) version instead of a hard
        label-flip count. \textbf{Trade-off:} cheap, continuous with the prior work, and it might
        clean up the outliers. But it doubles down on an axis the base paper's own nuance already
        suspects. If consistency is the wrong \emph{kind} of quantity, no amount of polishing fixes
        the outliers; you will have built a better thermometer for the wrong temperature.
  \item \emph{Bet on a single geometric quantity.} Drop consistency and pick \emph{one} geometric
        scalar instead: a bare margin (distance to the boundary), a bare sensitivity (a Lipschitz or
        gradient-norm estimate), or a \emph{margin-to-sensitivity ratio} that trades the two against
        each other. Part~I gives you reason to take the ratio family seriously --- the
        Lipschitz--margin certificate lower-bounds the robust radius by a margin over a Lipschitz
        constant --- but the certificate is background, not a verdict: it does not tell you that a
        ratio \emph{orders} models better than a bare margin or a bare sensitivity does, nor in which
        feature space to compute any of them. \textbf{Trade-off:} principled, and each candidate
        ties to geometry you can compute. But a margin and a Lipschitz constant are both nontrivial
        to measure on a real network; the certificate inequality is one-directional (a lower bound,
        not an equality --- you will revisit this in Q3); and you must decide \emph{which} feature
        space (raw pixels? the space the invariance leaves behind?) to measure in. Picking the wrong
        family, or the wrong space, is the way this direction fails.
  \item \emph{Let the data choose by regression.} Throw a basket of candidate scalars (consistency,
        margin, Lipschitz, gradient norms, boundary distance) at a panel of models and report which
        correlates best with measured robustness. \textbf{Trade-off:} maximally empirical and
        honest about what predicts what; it cannot be accused of theory-bias. But correlation on a
        model zoo confounds everything with everything (capacity, normalization, training), and a
        leaderboard of correlations is not an explanation --- it tells you \emph{that} something
        orders robustness, never \emph{why}.
\end{itemize}

\begin{problem}[commit before Part III]
First decide whether you even believe a single scalar can order robustness across architectures; if
not, say what minimal set of scalars you would need instead. If you do, pick the family you would bet
on --- consistency, a margin, a sensitivity, or a margin-to-sensitivity ratio --- and name the single
most likely reason your bet fails. If you chose a ratio or any geometric scalar, write down which
feature space you would measure it in and why raw pixels might be the wrong choice. If you chose
regression, write down the one confound you would control first.
\end{problem}

%----------------------------------------------------------------------
\subsection{Q2 --- What discriminative signal survives an imposed invariance?}
\label{sub:p2-q2}

\paragraph{Situation.} The linear story is brutally clean: force shift-invariance on a linear model
and it sees only the mean. But real networks are not linear classifiers on raw pixels; they push
data through nonlinear feature maps first. If a convolution followed by a squaring nonlinearity and
a pooling step is invariant to shifts, it is \emph{not} reduced to the mean --- it keeps something
richer. You want to know \emph{exactly what}. Not ``it keeps more''; the precise functional. Because
whatever survives the invariance is the only material left to build a margin out of, and Q1 told you
margin is the thing that matters.

\paragraph{Candidate directions.}
\begin{itemize}[leftmargin=1.4em]
  \item \emph{Quotient and diagonalize.} Treat the invariance as a projection: the model can only
        depend on the part of the input that the group leaves fixed, so write that projector down and
        diagonalize the group action (Fourier, for cyclic shifts) to read off the surviving
        coordinates in closed form. \textbf{Trade-off:} this is exactly the machinery Part~I built,
        it gives an \emph{exact} answer for the linear and quadratic cases, and it makes the margin
        computable. But it is cleanest for a single, idealized invariant layer; pushing it through a
        deep stack of real layers is where the rigor frays.
  \item \emph{Climb the ladder of classical invariants.} There is a known hierarchy: the mean
        (degree~1), the power spectrum / autocorrelation (degree~2, what a square-and-pool computes),
        the bispectrum (degree~3), and beyond. Ask which rung a given architecture sits on and what
        it can therefore separate. \textbf{Trade-off:} it connects the architecture to a beautiful,
        well-understood body of signal-processing theory, and it predicts \emph{failure} cases ---
        classes that one rung cannot tell apart but the next can. But higher rungs cost more
        (a larger Lipschitz constant; recall that any margin-to-sensitivity ratio has a sensitivity
        in its denominator), and you must show the
        architecture \emph{actually} realizes the rung, not just that the rung exists.
  \item \emph{Probe empirically.} Skip the analysis: train the invariant model and read out which
        input perturbations it is blind to, mapping the surviving signal by experiment.
        \textbf{Trade-off:} works for any architecture, no idealization. But it gives you a point
        cloud, not a formula; it will not tell you the \emph{margin} the surviving signal supports,
        and it cannot certify anything.
\end{itemize}

\paragraph{A pointed sub-question.} Recall the base paper's worst case: the ``one hot pixel''
example, $\pm e_j$, where the class is encoded by \emph{where} the spike sits. Two such inputs that
differ only by a shift have the \emph{same} power spectrum --- the degree-2 invariant is blind to
location. So if your architecture stops at degree~2, it cannot separate them at all, and the forced
trade-off bites hardest exactly there. Does a richer invariant rescue this corner, and at what cost
to the Lipschitz denominator? Hold that thought; it is a live fork.

\begin{problem}[commit before Part III]
For a convolution--square--pool layer, write your best guess at the surviving signal in one line.
Then decide: would you try to \emph{defeat} the hot-pixel worst case with a higher-order invariant,
or would you accept it as a genuine limit of low-degree invariance? State the price of your choice in
terms of the margin and the sensitivity (Lipschitz) it would move.
\end{problem}

%----------------------------------------------------------------------
\subsection{Q3 --- Does the surviving signal certify an $\ell_p$ radius?}
\label{sub:p2-q3}

\paragraph{Situation.} Suppose Q2 hands you an exact margin $\eta$ for the surviving signal, and
suppose you are pursuing the margin-to-sensitivity ratio of Q1's second direction, so the object on
the table is the Lipschitz--margin certificate $r_2 \ge \eta/L$ from Part~I. It is a \emph{lower}
bound: it guarantees no adversary inside radius $\eta/L$, but it says nothing about how much
\emph{further} the true safe radius might extend. For this to be the missing link the whole field
lacks --- a result tying translation-invariance to an actual $\ell_p$ radius --- you have to decide
what kind of statement you are even after. (Whether this ratio actually \emph{orders} models better
than its rivals is still Q1's open business; here you are only asking what the certificate, taken on
its own terms, can and cannot promise.)

\paragraph{Candidate directions.}
\begin{itemize}[leftmargin=1.4em]
  \item \emph{Settle for a one-sided certificate (and predictor).} Accept $r_2 \ge \eta/L$ as a
        guarantee and an empirical predictor, and do not claim it pins the radius exactly.
        \textbf{Trade-off:} honest and immediately usable; it is enough to \emph{order} models and to
        \emph{guarantee} safety up to the bound. But a lower bound alone cannot tell you whether a
        model with a small certificate is actually fragile or merely has a loose bound --- and a
        referee will ask whether your quantity \emph{characterizes} robustness or merely lower-bounds
        it.
  \item \emph{Chase a converse.} Try to prove an upper bound too, so the ratio brackets the true
        radius from both sides. \textbf{Trade-off:} a two-sided bracket is far stronger and would
        make the quantity a genuine measure of robustness, not just a safety floor. But you should
        first ask whether a \emph{universal} converse can even exist --- look for a degenerate example
        where the ratio is tiny yet the true radius is large. If you can build one, an unconditional
        converse is dead, and the honest move is to find the extra \emph{condition} under which a
        converse holds. (This is the search for what has no converse, a research move in its own
        right; see Section~\ref{sec:p2-taste}.)
  \item \emph{Add a second, invariance-specific radius.} The certificate above is the
        \emph{additive-noise} radius. But Tramer's excessive-invariance attack changes the true label
        \emph{along the orbit} while the prediction stays fixed --- a different failure the additive
        radius does not see. You could certify a separate ``orbit-flip'' radius for that mode.
        \textbf{Trade-off:} it captures a real and distinct threat and makes the ``invariance can
        also \emph{hurt}'' direction precise. But it complicates the story --- now there are two
        radii to report, and you must say when each one is the binding constraint.
\end{itemize}

\begin{problem}[commit before Part III]
Decide whether you would chase a converse or settle for a one-sided certificate. If you would chase
it, sketch the degenerate example you would build \emph{first} to check whether an unconditional
converse is even possible --- a function whose certificate is loose by an unbounded factor. If you
would not, say what claim you are giving up.
\end{problem}

%----------------------------------------------------------------------
\subsection{Q4 --- Does any of this survive a strong attack at matched capacity?}
\label{sub:p2-q4}

\paragraph{Situation.} Acts~I--IV are littered with the same two methodological landmines. First,
PGD/FGSM attacks can \emph{appear} to show robustness that vanishes under a stronger evaluation
because the gradients were masked, not the model genuinely robust (Part~I). The ``invariance helps''
evidence is almost entirely PGD/FGSM. Second, the architectures being compared differ in
\emph{capacity} (and in normalization, per \citet{galloway2019}), so any robustness gap might be a
capacity or normalization gap wearing an invariance costume. You want to know whether any geometric
story you settle on in Q1--Q3 is \emph{real} or an artifact of weak attacks and unmatched models.

\paragraph{Candidate directions.}
\begin{itemize}[leftmargin=1.4em]
  \item \emph{Trust a stronger attack ensemble.} Re-evaluate every model with a standardized,
        parameter-free, ensemble attack designed precisely to catch gradient masking, and treat that
        number as ground truth. \textbf{Trade-off:} this is the community standard and the only way to
        kill the gradient-masking objection. But it is expensive, and a single robust-accuracy number
        at one perturbation budget hides the \emph{radius} you actually care about.
  \item \emph{Measure the radius directly with a min-norm attack.} Instead of ``what fraction
        survives at budget $\epsilon$,'' find, per input, the \emph{smallest} perturbation that flips
        it --- the empirical robust radius. \textbf{Trade-off:} this is the right quantity to compare
        against an $\eta/L$ \emph{radius} certificate, and it gives a distribution, not one number.
        But min-norm attacks are themselves attacks and can be fooled; you should cross-check them
        against the ensemble above.
  \item \emph{Control capacity by construction.} Build a family of architectures --- a plain one, an
        anti-aliased one, an exactly-cyclic-invariant one, a shift-augmented one --- with
        \emph{identical} parameter counts, the normalization confound pinned down, and the
        \emph{same} training. Then any robustness difference is attributable to the invariance, not
        to size. \textbf{Trade-off:} this is the strongest possible causal isolation and lets you ask
        not ``does invariance help?'' but ``which term, margin or Lipschitz, does \emph{this}
        invariance move?'' But matching capacity across structurally different models is fiddly, and
        you must verify the matched models really are comparable in every nuisance dimension.
  \item \emph{Stay on synthetic data where the truth is known.} Run everything on the constructions
        where the margin and radius are computable in closed form, sidestepping attack noise
        entirely. \textbf{Trade-off:} clean and unimpeachable, and it isolates mechanism. But a
        referee --- and the gap from Act~IV --- will say the open question is precisely about
        \emph{real} data under a \emph{strong} attack; synthetic-only re-opens the very gap you set
        out to close.
\end{itemize}

\begin{problem}[commit before Part III]
Design your headline experiment in three lines: which models you compare, what you hold fixed to
make the comparison causal, and which attack(s) you trust for the final number. Name the one
confound that, if a reviewer raised it, would most damage your conclusion --- and say how your design
already answers it.
\end{problem}

%----------------------------------------------------------------------
\subsection{Q5 --- Why do trained models land where the diagnostic put them?}
\label{sub:p2-q5}

\paragraph{Situation.} Q1--Q4 are about what a \emph{given} model's geometry implies. But the
Act~II quadrant diagram was about \emph{trained} models: gradient descent, not a fixed weight
vector, decides where each architecture ends up on the consistency-versus-robustness plane. So even
a perfect geometric theory leaves a gap: \emph{why does optimization land an architecture in the
quadrant it lands in?} Answer this and the heuristic quadrant diagram becomes a principle; leave it
and you have explained the map but not the territory. Part~I's implicit-bias and NTK lenses are the
tools here.

\paragraph{Candidate directions.}
\begin{itemize}[leftmargin=1.4em]
  \item \emph{Reduce to a regime where the bias is known.} In the lazy/NTK regime the trained
        network behaves like a fixed kernel method, where ``what gradient descent selects'' has a
        clean answer. Ask what enforcing invariance does to that kernel's solution. \textbf{Trade-
        off:} tractable, and it can turn an empirical trend into a theorem about the selected
        solution. But the lazy regime is an idealization (wide, near-linear training), the activation
        matters (a smooth activation behaves very differently from bare ReLU under this lens), and
        what holds lazily may not hold in the feature-learning regime real training lives in.
  \item \emph{Split by data geometry.} The Act~III/IV papers suggest the answer is conditional. On
        data whose classes are \emph{closed under the symmetry} (every shift of a class member is in
        the class), enforcing invariance may cost nothing --- gradient descent might symmetrize for
        free. On near-orthogonal \emph{cluster} geometry (the Frei/Li/Min-Vidal setting) invariance
        may genuinely help the ratio. Prove the easy case, conjecture the hard one.
        \textbf{Trade-off:} this matches the contradiction's actual shape --- different geometries,
        different verdicts --- and gives you a provable positive result plus an honest open problem.
        But it splits the question into pieces, one of which you will have to leave open, and you must
        be candid that the cluster case resists a full trajectory-level proof.
  \item \emph{Stay empirical and characterize the landing pattern.} Train the matched families from
        Q4 across seeds and data regimes and report \emph{where} they land and how tightly,
        upgrading the quadrant diagram into a measured phenomenon without proving it.
        \textbf{Trade-off:} robust and honest, and it documents exactly the regime-dependence the
        theory should later explain. But ``we observed where models land'' is a description, not the
        principle the Act~II diagram was missing; a reviewer asks \emph{why}.
\end{itemize}

\begin{problem}[commit before Part III]
Decide whether you would attack this in the lazy regime, split it by data geometry, or characterize
it empirically --- and say which part of the question you are willing to leave \emph{open}. (Honest
answer: at least one piece of Q5 is still open in the project. Predict which piece, and why it is the
hard one.)
\end{problem}

%======================================================================
\section{Research taste: how to find questions like these}
\label{sec:p2-taste}

The five questions above did not arrive labelled. Someone had to manufacture them from a tangle of
contradictory papers and an unexplained scatter plot. This section is about that manufacturing.
It is not a checklist --- taste is not a checklist --- but there are recurring \emph{moves} that
generate good questions, and there are real \emph{criteria} for deciding which questions are worth
the months they cost. I illustrate each move with a situation from this project, without giving the
answer, so you can practice the move rather than memorize the result.

\subsection{Recurring moves}
\label{sub:p2-moves}

\paragraph{Quotient out the nuisance.} When a symmetry or a degree of freedom is making the problem
look complicated, do not fight it --- divide by it. Ask what is left once the nuisance is collapsed.
\emph{In this project:} ``the model is shift-invariant'' looks like a statement about infinitely many
shifted inputs. Quotient by the shift group and it becomes a statement about a single object --- the
projection of the data onto the part the group fixes. The infinitely-many-shifts problem becomes a
one-projector problem. (This is the engine behind Q2.)

\paragraph{Diagonalize the symmetry.} A symmetry that commutes with your operator can be diagonalized,
and in the eigenbasis the operator is just multiplication --- everything decouples. \emph{In this
project:} circular shifts are diagonalized by the Fourier basis, so an invariant's action on the data
becomes coordinate-wise on frequencies, and ``what survives invariance'' becomes ``which frequency
coordinates the invariant can still read.'' The hard algebra of convolutions turns into reading off a
diagonal. Whenever you see a group acting linearly, ask for its eigenbasis first.

\paragraph{Bracket what has no converse.} A one-sided inequality is a confession that you do not yet
know the other side. Before assuming both sides hold, deliberately hunt for a counterexample to the
converse. If you find one, you have learned that the clean two-sided statement is \emph{false in
general} --- which is itself a result --- and the real question becomes ``under what extra condition
does a converse hold?'' \emph{In this project:} the certificate $r_2 \ge \eta/L$ is one-sided; the
right reflex is to ask whether a universal converse can exist before trying to prove one. (This is
the heart of Q3, and the degenerate-example reflex below.)

\paragraph{Find the degenerate case.} The fastest way to test a claim or kill a converse is to push
it to an extreme where one term blows up or collapses, and see if the claim survives. \emph{In this
project:} to probe whether $\eta/L$ pins the radius, you reach for a function whose curvature can be
cranked arbitrarily high so the certificate becomes arbitrarily loose while the true radius does not
move. A single such example settles a general question instantly. Degenerate cases are cheap and
decisive; build them \emph{before} you invest in a hard proof.

\paragraph{Reduce to the per-unit atom.} A network is a composition of many small pieces. Instead of
reasoning about the whole, find the smallest piece that carries the phenomenon and solve \emph{that}
exactly, then ask how the pieces compose. \emph{In this project:} rather than analyze a deep
invariant network, isolate one convolution--nonlinearity--pool unit and compute exactly what it keeps
and what margin it supports. The atom is solvable; the composition is the next question. If you
cannot solve the atom, you certainly cannot solve the network --- and if you can, you have a building
block and a sanity check for everything larger.

\paragraph{Transport a finite lemma to function space.} A clean fact about finite-dimensional vectors
(projections are non-expansive; averaging cannot increase a norm) often has an exact analogue in a
function space (a reproducing-kernel Hilbert space) that turns an empirical trend about trained
networks into a theorem. \emph{In this project:} ``averaging over the orbit cannot increase the norm''
is a one-line fact for vectors; transported to the kernel that describes a wide network's training, it
becomes a statement about the solution gradient descent selects. The move is: prove it where it is
trivial, then carry it across the bridge that your regime (e.g.\ the lazy/NTK lens from Part~I)
provides. The catch is always whether the bridge holds --- which is exactly why Q5's lazy direction
carries an activation caveat.

\subsection{Which questions are worth asking}
\label{sub:p2-worth}

Moves generate candidate questions faster than you can answer them. Triage is the other half of
taste. Three criteria did the real work here.

\paragraph{1. There is a small, testable instance.} A question you can probe on a two-frequency toy
or a closed-form linear case \emph{this afternoon} is worth a hundred you can only settle with a
six-month training run. The small instance tells you fast whether the big claim is even plausible,
and it doubles as a sanity check on the eventual general proof. Every one of Q1--Q5 has a synthetic
miniature where the answer is computable; that is not a coincidence, it is the selection filter.
If you cannot construct a small instance, you usually do not yet understand the question well enough
to ask it.

\paragraph{2. It sits exactly at a literature contradiction.} The most fertile questions live on the
seam where two credible papers disagree, because the disagreement guarantees that \emph{someone}
is wrong about \emph{something specific}, and resolving it cannot be a no-op. Q1 sits on the seam
between ``consistency orders robustness'' and Ge's own ``margin predicts better than dimension.''
Q5 sits on the seam between ``invariance hurts'' (Frei/Li) and ``invariance helps'' (Saha--Gokhale,
Chen--Zhu). A question that no published paper would dispute is rarely worth asking; a question that
forces two published papers to reconcile almost always is.

\paragraph{3. The answer changes a decision.} Ask what someone would \emph{do} differently once they
know the answer. Q1's answer changes which scalar you compute to screen architectures. Q4's answer
changes whether you trust a PGD number at all. Q3's answer changes whether you can advertise a
\emph{certificate} or only a \emph{predictor}. If both possible answers lead to the same action, the
question is decorative; drop it.

\paragraph{Be honest: some good-looking questions did not pan out.} Taste includes being wrong
gracefully. Several attractive questions in this project bought less than they promised. Polishing the
consistency metric (Q1's first direction) was tempting and continuous with prior work, but if the
axis is the wrong \emph{kind} of quantity, a better thermometer does not help --- a lesson you only
learn by trying it or by reasoning hard about why not. Chasing an \emph{unconditional} converse for
the certificate (Q3) is a dead end the degenerate-case move kills in an afternoon; the salvageable
question is the \emph{conditional} one, which is strictly humbler. And the cleanest part of Q5 ---
the optimization story on symmetric data --- yields a real theorem, while the part everyone wants ---
the cluster-geometry case in the realistic feature-learning regime --- stays a conjecture supported
by experiments, not a proof. A primer that pretended otherwise would be teaching you the wrong
lesson. Knowing which of your questions is provable, which is merely true, and which is only plausible
is not a footnote to the research. It \emph{is} the research.

\medskip

You now have the five questions, the moves that generate them, and the criteria that select them.
Before you read Part~III, go back to your five committed answers from the \textit{commit before
Part~III} prompts. Write, for each, the single observation that would most change your mind. Then,
and only then, turn the page.

% =====================================================================
%  CITATION KEYS USED IN THIS FRAGMENT
% =====================================================================
%  REUSED (already in primer.bib):
%    li2025feature   -- Li et al. "Feature averaging..." (the Li 2024 paper;
%                       primer.bib already keys it as li2025feature)
%
%  NEW KEYS -- must be added to primer.bib (real, verified references):
%    ge2021shift      -- Ge, Kuang & Zou (or as published), "Shift Invariance
%                        Can Reduce Adversarial Robustness", NeurIPS 2021.
%    kamath2020       -- Kamath, Deshpande, Subrahmanyam (et al.),
%                        "Invariance vs. Robustness of Neural Networks",
%                        ICLR 2020 / OpenReview 2020.
%    kamath2021       -- Kamath et al., "Can we have it all? On the trade-off
%                        between spatial and adversarial robustness ...",
%                        NeurIPS 2021 (introduces CuSP).
%    galloway2019     -- Galloway, Golubeva, Tanay, Moussa, Taylor,
%                        "Batch Normalization is a Cause of Adversarial
%                        Vulnerability", ICML 2019 workshop / arXiv:1905.02161.
%    zhang2019blurpool-- R. Zhang, "Making Convolutional Networks
%                        Shift-Invariant Again", ICML 2019 (BlurPool/anti-alias).
%    tramer2020       -- Tramer, Behrmann, Carlini, Papernot, Jacobsen,
%                        "Fundamental Tradeoffs between Invariance and
%                        Sensitivity to Adversarial Perturbations", ICML 2020.
%    frei2023         -- Frei, Vardi, Bartlett, Srebro, "The double-edged
%                        sword of implicit bias: generalization vs.
%                        robustness in ReLU networks", NeurIPS 2023.
%    minvidal2024     -- Min & Vidal, "Can Implicit Bias Imply Adversarial
%                        Robustness?", ICML 2024 (polynomial-ReLU flips bias).
%    sahagokhale2024  -- Saha & Gokhale, "Translation-Invariant Polyphase
%                        Sampling (TIPS)", WACV 2024.  NOTE: the local file is
%                        mislabeled "rajput2024_..."; correct authorship is
%                        Saha & Gokhale.
%    wang2025bridging -- (Wang et al.) "Bridging Symmetry and Robustness:
%                        ... Equivariance ...", 2025 (rotation/scale equivariance).
%    chenzhu2023      -- Chen & Zhu, "Implicit Bias of (Linear) Equivariant
%                        Steerable Networks", 2023.
%    lawrence2021     -- Lawrence et al., "Implicit Bias of Linear Equivariant
%                        Networks", 2021.
% =====================================================================
